\section*{参考文献}
\phantomsection
\addcontentsline{toc}{section}{参考文献}

\begin{enumerate}
\def\labelenumi{(\Roman{enumi})}
\item
  O. NEUGEBAUER，《古代数学科学史讲义》，第 I 卷：希腊以前的数学，柏林（Springer），1934。
\item
  J. TROPFKE，《初等数学史》，第 IV–VI 卷，柏林–莱比锡（de Gruyter），1923–24。
\item
  欧几里得《几何原本》，5 卷，J. L. Heiberg 编，莱比锡（Teubner），1883–88。
\end{enumerate}

(III bis) T. L. HEATH，《欧几里得〈几何原本〉十三卷》\ldots，3 卷，剑桥，1908。

\begin{enumerate}
\def\labelenumi{(\Roman{enumi})}
\setcounter{enumi}{3}
\item
  T. L. HEATH，《佩尔加的阿波罗尼奥斯：圆锥曲线论》，剑桥（Univ. Press），1896。
\item
  Ptolemaei Cl. Opera（托勒密全集），J. L. Heiberg 编，2 卷，莱比锡（Teubner），1898–1903。
\item
  G. DESARGUES，《全集》\ldots 第 I 卷，巴黎（Leiber），1864：Brouillon project d'une atteinte aux événements des rencontres d'un cône avec un plan，第 103-230 页。
\item
  P. FERMAT，《全集》，第 I 卷，巴黎（Gauthier-Villars），1891：a) Ad locos planos et solidos Isagoge，第 91-110 页（法文译文见同书第 III 卷，第 84-101 页）；b) Isagoge ad locos ad superficiem，第 111-117 页（法文译文见同书第 III 卷，第 102-108 页）。
\item
  L. EULER：a) Introductio in Analysin Infinitorum（Opera Omnia (1)，第 IX 卷，苏黎世–莱比锡–柏林（O. Füssli 与 B. G. Teubner），1945）；b) Theoria motus corporum solidorum seu rigidorum（Opera Omnia (2)，第 III 卷，苏黎世–莱比锡–柏林（O. Füssli 与 B. G. Teubner），1948）；c) Problema algebraicum ob affections prorsus singulares memorabile（Opera Omnia，(1)，第 VI 卷，莱比锡–柏林（Teubner），1921，第 287-315 页）；d) Formulae generales pro translatione quacunque corporum rigidorum，Novi Comm. Acad. Sc. imp. Petrop.，第 XX 卷（1776），第 189-207 页；e) De centro similitudinis.（Opera Omnia (1)，第 XXVI 卷，苏黎世（O. Füssli），1956，第 276-285 页）。
\item
  J. L. LAGRANGE，《全集》，巴黎（Gauthier-Villars），1867-1892：a) 《关于极大极小方法的研究》，第 I 卷，第 3–20 页；b) 《算术研究》，第 III 卷，第 695–795 页。
\item
  G. MONGE，《画法几何》，巴黎，1798。
\item
  C. F. GAUSS，《全集》：a) Disquisitiones arithmeticae，第 I 卷，哥廷根，1870；b) 《空间的变换》，第 VIII 卷，哥廷根，1900，第 357–362 页。
\item
  J.-V. PONCELET，《图形的射影性质论》，第 I 卷，第 2 版，巴黎（Gauthier-Villars），1865。
\item
  A. F. Möbius，《全集》，莱比锡（Hirzel），1885–87：a) 《重心演算》，第 I 卷，第 1–388 页；b) 《论空间图形之间一类特殊的对偶关系》，第 I 卷，第 489–515 页（= Crelle's Journal，第 X 卷，1833）；c) 《论解析球面几何的一种新处理》，第 II 卷，第 1–54 页。
\item
  A.-L. CAUCHY：a) 《关于无穷小演算在几何中的应用的讲义》（《全集》，(2)，第 V 卷，巴黎（Gauthier-Villars），1903）；b) 《论用以确定行星长期摄动的方程》（《全集》，(2)，第 IX 卷，巴黎（Gauthier-Villars），1891，第 174–195 页）。
\item
  M. CHASLES：a) 《关于两个物体所成系统的一般性质的注记》，Bull. de Férussac，第 XIV 卷（1830），第 321–326 页；b) 《几何方法的起源与发展的历史概述》，布鲁塞尔，1837。
\item
  E. GALOIS，《数学著作》，巴黎（Gauthier-Villars），1897。
\item
  W. R. HAMILTON，《四元数讲义》，都柏林，1853。
\item
  A. CAYLEY，《数学论文集》，剑桥，1889-1898：a) 《关于与四元数有关的若干结果》，第 I 卷，第 123-126 页（= Phil. Mag.，1845）；b) Sur les déterminants gauches，第 I 卷，第 410-413 页（= J. de Crelle，第 XXXVIII 卷（1848））；c) Recherches ultérieures sur les déterminants gauches，第 II 卷，第 202-215 页（= J. de Crelle，第 L 卷（1855））；d) 《关于代数型的第六篇论文》，第 II 卷，第 561-592 页（= Phil. Trans.，1859）。
\item
  C. G. J. JACOBI，Gesammelte Werke（《全集》），柏林（G. Reimer），1881-1891：a) 《论普法夫方法》\ldots 第 IV 卷，第 17–29 页；b) 《论一个基本的代数定理及其应用》，第 III 卷，第 593–598 页。
\item
  J. J. SYLVESTER，《数学论文集》，第 I 卷，剑桥，1904：《每个齐次二次多项式均可经实正交代换化为正负平方之和的形式这一定理的证明》，第 378–381 页（= Phil. Mag.，1852）。
\item
  E. LAGUERRE，《全集》，第 II 卷，巴黎（Gauthier-Villars），1905。
\item
  C. HERMITE，《全集》，第 I 卷，巴黎（Gauthier-Villars），1905：《论二次型理论》，第 200–263 页（= J. de Crelle，第 XLVII 卷（1854））。
\item
  K. G. V. von STÄUDT，《位置几何学的贡献》，纽伦堡，1856。
\item
  E. BELTRAMI：a) 《非欧几何解释试论》，Giorn. di Mat.，第 VI 卷（1868），第 284–312 页；b) 《常曲率空间的基本理论》，Ann. di Mat. (2)，第 II 卷（1868–69），第 232–255 页。
\item
  F. KLEIN，《数学文集》，第 I 卷，柏林（Springer），1921：a) 《论所谓非欧几何》，第 254-305 页（= Math. Ann.，第 IV 卷（1871））；b) 《关于近代几何研究的比较考察》，第 460-497 页（= Math. Ann.，第 XLIII 卷（1893））。
\item
  C. JORDAN，《置换与代数方程论》，巴黎（Gauthier-Villars），1870。
\item
  G. FROBENIUS：a) 《论线性置换与双线性型》，J. de Crelle，第 LXXXIV 卷（1878），第 1-63 页；b) 《整系数线性型的理论》，J. de Crelle，第 LXXXVI 卷（1879），第 146-208 页。
\item
  W. K. CLIFFORD，《数学论文集》，伦敦（Macmillan），1882：a) 《论几何代数的分类》，第 397-401 页；b) 《格拉斯曼广延代数的应用》，第 266-276 页（= Amer. Journ. of Math.，第 I 卷（1878））。
\item
  R. LIPSCHITZ，《关于平方和的研究》，波恩，1886。
\item
  D. HILBERT，《几何学基础》，莱比锡（Teubner），1899。
\item
  E. WITT，《任意体上的二次型理论》，J. de Crelle，第 CLXXVI 卷（1937），第 31-44 页。
\item
  E. CARTAN，《旋量理论讲义》，Actual. Sci. et Industr.，第 643 与 701 号，巴黎（Hermann），1938。
\item
  C. L. SIEGEL，《辛几何》，Amer. Journ. of Math.，第 LXV 卷（1943），第 1-86 页。
\item
  C. CHEVALLEY，《旋量的代数理论》，纽约（Columbia Univ. Press），1954。
\item
  M. EICHLER，《二次型与正交群》，柏林–哥廷根–海德堡（Springer），1952。
\end{enumerate}

